\documentclass[11pt]{article}

\usepackage[margin=1in]{geometry}
\usepackage{amsmath,amssymb,amsthm}
\usepackage{graphicx}
\usepackage{booktabs}
\usepackage{xcolor}
\usepackage{hyperref}
\hypersetup{colorlinks=true, linkcolor=blue!60!black, citecolor=blue!60!black, urlcolor=blue!60!black}

\newtheorem{theorem}{Theorem}
\newtheorem{corollary}{Corollary}

\newcommand{\peff}{p_{\mathrm{eff}}}
\newcommand{\conv}{\mathrm{conv}}
\newcommand{\rhoalign}{\rho_c^{\mathrm{align}}}
\newcommand{\rhoperc}{\rho_c^{\mathrm{perc}}}
\newcommand{\Pperc}{P_{\mathrm{perc}}}
% Figure helper: include if present, framed placeholder otherwise (repo figures)
\newcommand{\repofig}[2]{\IfFileExists{#1}{\includegraphics[width=#2]{#1}}{\fbox{\parbox{#2}{\centering\vspace{2em}\texttt{\detokenize{#1}}\\ (figure in \texttt{figures/paper3/})\vspace{2em}}}}}

\title{Percolation of the arrow of time:\\
an exact factorization law for lattices of shared quantum environments}
\author{Mostfa Lejri\\ \small Independent Researcher}
\date{\small Draft: July 2026}

\begin{document}
\maketitle

\begin{abstract}
In companion papers we showed that decoherence fixes a local arrow of time
equal to the record information of a monitored T-odd observable, and that
competing arrows are arbitrated by the preparation of shared records. Here we
ask how many local arrows knit into one. We study lattices of chiral rings,
each monitored by a private recording bath and sharing witnesses with its
neighbours, with a density $\rho$ of contrarian rings prepared with
record-consuming (time-mirrored) private baths. Three results follow. First,
a single universal conversion curve $\conv(f)$, the probability that sharing
fraction $f$ converts a contrarian, is independent of $\rho$ (spread $0.007$)
and carries all the model physics. Second, an exact factorization law: under fresh
shared witnesses the arrow field reduces to independent-site occupation
$\peff = 1-\rho\,(1-\conv(f))$, with bulk conversion correlations compatible
with zero ($-0.025$) and a single finite-range nearest-neighbour correction
($+0.11$, from shared links) that leaves the universality class unchanged.
Both thresholds then follow in closed form with no free parameter: global
reversal at $\rhoalign = 1/[2(1-\conv)]$ (residual $0.003$) and percolation
of the aligned domain at $\rhoperc = (1-p_c)/(1-\conv)$ (mean residual
$0.022$). The universality class follows from the reduction itself: a
synthetic independent-site field is statistically indistinguishable from the
quantum lattice (mean deviation $0.040$, below the $0.056$ sampling floor),
so the aligned-domain transition inherits standard 2D percolation. Direct
exponent estimates are consistent ($\gamma/\nu = 1.78$ against $1.79$, quasi
exact and robust; $\nu_{\mathrm{eff}}$ drifting toward $4/3$ with size) and
are presented as consistency checks, not measurements. Third, the
factorization is a property of the boundary condition: charging the shared
witnesses (prescription M) reverses the flow. A majority of fresh rings is
converted backward and the global alignment flips sign, the F/M duality of
the companion paper operating at the scale of the whole lattice. The
direction of global time, in this model, is decided by the preparation of
the records the world shares, compressed into one measurable curve.
\end{abstract}

\tableofcontents

\section{Introduction}

Two companion papers built, in an exactly solvable model, a quantitative
bridge between the fixation of a time-reversal-odd observable and the
irreversibility of the monitoring process. The first \cite{p1} showed that
for a single three-site chiral ring monitored by recording qubits the
unconditional stochastic entropy production vanishes identically (the arrow
of time is spontaneously broken, existing only relative to a branch) and
that the branch-relative entropy production equals exactly the accumulated
distinguishability $D(t)$ of the environmental records. The second \cite{p2}
shared the environment between rings: a common bath fixes only what it can
distinguish, joint irreversibility carries an irreducibly collective
component with the exact decomposition $\Delta = I + C_{\mathrm{bwd}}$, and
when two rings with opposed arrows compete for shared witnesses, exactly one
arrow flips, at a sharp threshold $f^* \approx 0.40 < 1/2$ whose position,
and the identity of the loser, is set by the preparation of the shared
records (the F/M duality). Both papers end on the same conjecture, stated in
\cite{p1} as a contagion of local arrows into a consistent global one
(records begetting records, in the Darwinian sense of
\cite{zurek2009,riedel2012}) and in \cite{p2} as the natural next step: a
lattice of rings with nearest-neighbour bath overlap, asking whether
alignment percolates.

This paper answers that question, and the answer has more structure than the
conjecture asked for. The many-arrow problem does not merely resemble a
percolation problem; under the fresh shared-witness prescription it
\emph{reduces} to one. Three results organize the paper, and we announce
their epistemic statuses with them, continuing the practice of the
companions. First (measured, with its universality mechanically explained):
a single conversion curve $\conv(f)$, the probability that a contrarian ring
is converted by sharing fraction $f$, is independent of the contrarian
density $\rho$ and carries all the model physics
(Sec.~\ref{sec:conversion}). Second (the factorization law, carrying the paper; two of its three
structural facts are proven, the third verified as a positivity identity,
and a three-step numerical falsification programme fails to break it): the
arrow field
is an independent site-percolation field of occupation
$\peff = 1 - \rho(1-\conv(f))$, every collective threshold follows in closed
form with no free parameter, and the universality class of the
aligned-domain transition is inherited from standard two-dimensional
percolation by equivalence rather than measured
(Secs.~\ref{sec:factorization} and \ref{sec:global}). Third (the boundary,
exact in its mechanism): the reduction is a property of the record boundary
condition. Charging the shared witnesses reverses the global arrow and
destroys the factorization itself, the F/M duality of \cite{p2} operating at
the scale of the whole lattice (Sec.~\ref{sec:M}).

Section~\ref{sec:model} defines the lattice model, the local arrow, and the
validation gates. Section~\ref{sec:conversion} establishes the exact
baseline, the one-dimensional crossover, and the universal conversion curve.
Section~\ref{sec:factorization} states the factorization law and its
numerical proof. Section~\ref{sec:global} treats the percolation of the
global arrow and the status of the exponents. Section~\ref{sec:M} sets the
boundary of the reduction. Section~\ref{sec:discussion} consolidates the
epistemic statuses and the openings.

\section{The lattice model}
\label{sec:model}

\subsection{Rings, shared bonds, and the sharing fraction}

Each site $i$ of a lattice (a chain of length $L$ in 1D, an $L \times L$
square lattice with periodic boundary conditions in 2D) carries the
three-site chiral ring of \cite{p1}: a conserved chirality $s_i = \pm 1$,
the T-odd pointer doublet. The joint branch is $S = (s_1,\dots,s_{N_r})$ for
$N_r$ rings, and the record variable of witness $n$ is the linear form
$a_n(S) = \sum_i s_i\, G_{i,n}$ of \cite{p2}, with $G_{i,n} \ge 0$ the
coupling matrix. Every ring owns a private bath of $N_p$ qubits, coupled to
it alone; every nearest-neighbour bond $\langle i,j\rangle$ carries $N_s$
shared qubits, coupled equally to both its rings: witnesses that record the
compound variable $s_i + s_j$, the lattice version of the shared witnesses
of \cite{p2}. The structural sharing fraction is the fraction of a ring's
accessible witnesses that are shared,
%
\begin{equation}
f \;=\; \frac{\deg \cdot N_s}{N_p + \deg \cdot N_s},
\label{eq:f}
\end{equation}
%
with $\deg$ the coordination number ($\deg = 4$ in the 2D bulk). This
convention is the one implemented by the 2D production runs (their $f$
grid reproduces Eq.~\eqref{eq:f} exactly at $N_p=8$); the 1D runs used a
per-bond label, converted in Sec.~\ref{sec:conversion}.
Throughout, $f$ is swept structurally, by the integer $N_s$ at fixed
coupling scale, never by the coupling strength: strong shared couplings fold
the record bias $|\sin|$ and pollute the arrow with recurrences, the lattice
version of the finite-size revivals removed by the collision protocol in
\cite{p1}. Parameters throughout: $N_p = 8$ (2D) or $3$ (1D), collision
window $\Delta t = 0.3$, charge time $T = 1.0$, couplings drawn as in
\cite{p1} with scale $0.25$, true branch $S_{\mathrm{true}} = (+1,\dots,+1)$.

\subsection{Contrarians and the collision protocol}

A fraction $\rho$ of rings, chosen uniformly at random (the contrarian
minority), is prepared with a mirror-charged private bath: following the
mirror preparation of \cite{p2}, their private witnesses start in the
$\Theta$-conjugate of a forward evolution of duration $T$, so their records
are consumed rather than accumulated and their branch-relative entropy
production is negative. The remaining fraction $1-\rho$ (the fresh majority)
owns record-free private baths. The shared witnesses on every bond are
prepared fresh, prescription F of \cite{p2}, except in Sec.~\ref{sec:M},
where they are charged (prescription M). All preparations enter through the
unified record bias of \cite{p2},
%
\begin{equation}
\beta_n(S) \;=\; \sin\!\big(2\, a_n(S)\, (\Delta t - T_c)\big),
\qquad T_c = 0 \;(\mathrm{F,\ fresh}), \quad T_c = T \;(\mathrm{M,\ charged}),
\label{eq:beta}
\end{equation}
%
and each witness is a collision-protocol record: it interacts during one
window, is measured once, and its record is permanent. Records therefore
accumulate monotonically, the ratchet regime of \cite{p1}, and the arrow of
the next subsection is well defined.

\subsection{The local arrow}
\label{sec:arrow}

For a single ring, \cite{p1} identified the branch-relative entropy
production with the record distinguishability $D$; for two competing rings,
\cite{p2} used the sign of the marginal production $\sigma_i$ as the arrow.
On the lattice we refine the state sign into a rate sign:
%
\begin{equation}
\tau_i \;=\; \mathrm{sign}\!\left(\frac{d\langle \sigma_i^{\mathrm{loc}}\rangle}{dt}\right),
\label{eq:tau}
\end{equation}
%
the direction of \emph{accumulation} of ring $i$'s accessible record
information: growing (forward, $\tau_i = +1$) or being consumed (backward,
$\tau_i = -1$). Here $\sigma_i^{\mathrm{loc}}$ is built from the marginal
likelihood ratio of \cite{p2} restricted to ring $i$'s accessible witnesses
(its private bath plus every shared bond touching it), with the neighbour
chiralities marginalized. Two remarks fix the status of this definition.

\emph{Why a rate, not a state sign.} The inference ratio
$\langle\sigma_i\rangle$ measures how much has been fixed, a stock; the
physically contagious quantity, what a shared witness transmits between
rings, is the direction in which that stock is currently moving. The two
signs can disagree only in a transient (a stock still positive but already
being consumed), which is precisely the window where non-permanent records
produce recurrences, hence the collision protocol. At $f = 0$ the two
definitions coincide exactly for the canonical preparations: a fresh ring
has $\langle\sigma_i\rangle = +D_i$ and growing, a charged ring
$\langle\sigma_i\rangle = -D_i$ and consuming, reproducing the C5 limit of
\cite{p2} at machine precision (gate G6 below). Paper 3 does not contradict
its companions; it sharpens the arrow from a state sign into a rate sign, in
exactly the regime (permanent records) where both are monotone.

\emph{Why the local marginal.} The exact global marginal costs $2^{N_r}$ and
is unusable beyond a handful of rings. The local marginal costs
$O(2^{\deg+1})$ per site, and gate G4 shows it agrees with the global
marginal in sign at every site tested (mean correlation $0.97$). The local
marginal is therefore the production definition: lattices of $256$ rings
($L=16$) are tractable, and the $2^{N_r}$ object is never built. It is also
the physically honest choice: $\tau_i$ is the arrow assigned by an observer
reading only ring $i$'s accessible records, the lattice version of the
operational marginalization of \cite{p2}.

\subsection{Validation gates}
\label{sec:gates}

Following the methodology of both companion papers, no physics run is
performed, and no figure produced, until an independent verification suite
passes. The six gates connect the lattice engine to every exact result of
the programme:

\begin{center}
\begin{tabular}{llr}
\toprule
Gate & Statement & Residual \\
\midrule
G1 & $\Delta = I + C_{\mathrm{bwd}}$ (\cite{p2}, Thm.~1) on the lattice-generated record law & $2.8\times 10^{-17}$ \\
G2 & $f=0$ reduction to \cite{p1}, Thm.~2 ($\langle\sigma_i\rangle = D_i$) & $2.2\times 10^{-16}$ \\
G3 & $L=3$ brute-force state vector vs.\ branch construction & $2.9\times 10^{-17}$ \\
G4 & local marginal $\equiv$ global marginal (sign $100\%$, corr.\ $0.97$) & sampling \\
G5 & collision protocol: exact additivity for disjoint baths, $\langle\sigma_i\rangle = D_i$ & $0.0$ \\
G6 & collision $L=2$ reproduces the \cite{p2} C5 limit at $f=0$ ($+D$, $-D$) & $2.2\times 10^{-16}$ \\
\bottomrule
\end{tabular}
\end{center}

\section{The universal conversion curve}
\label{sec:conversion}

\subsection{The exact baseline}

One exact test runs through every figure of the paper. At $f = 0$ no witness
is shared: fresh rings point forward, charged rings backward, and the signed
alignment must equal the discrete baseline
%
\begin{equation}
A(f{=}0) \;=\; 1 - 2\,\frac{n_c}{N_r}, \qquad n_c = \mathrm{round}(\rho N_r),
\label{eq:baseline}
\end{equation}
%
the integer contrarian count. It does, with deviation $0.0$ throughout
(exact). Comparing to the continuum $1-2\rho$ instead produces a spurious
deviation of order $1/L$, a discretization trap worth recording.

\subsection{One dimension: contagion, but a crossover}

On the chain (Figs.~\ref{fig:chain} and \ref{fig:domains}), the signed
alignment $A(f)$ rises from the baseline \eqref{eq:baseline} toward full
alignment as sharing grows. A convention note is required here, recorded
for reproducibility: the 1D production runs labelled their axis with the
per-bond fraction $f_b = N_s/(N_p + 2N_s)$, which at $\deg = 2$ is exactly
half the per-ring convention of Eq.~\eqref{eq:f}; all values quoted in this
paper are in the unified per-ring convention, $f = 2 f_b$. In that
convention the conversion onset sits at $f \approx 0.67$--$0.80$, the same
high-sharing neighbourhood as the 2D conversion curve of the next
subsection, and as the pair threshold under the same fresh-shared
preparation with a charge head start ($f^*_{\mathrm{F}} \approx 0.76$ at
$t = T/2$ in \cite{p2}); the equal-effort point $f^* \approx 0.40$ of
\cite{p2} is not the relevant comparison, since the lattice contrarian,
like the head-start mirror ring, begins fully charged. But the transition
is a crossover, not a phase transition (numerical): the mean
$\tau$-domain length grows with $f$ while the $L = 8/16/32$ curves
overlap, with no finite-size divergence, consistent with the absence of a
1D site-percolation transition ($p_c = 1$) \cite{stauffer}. One dimension
motivates two.
% TODO(P3-F, F3): regenerate the two chain figures with the per-ring axis
% (multiply the labelled f by 2) so figures and text share one convention.

\begin{figure}[t]
\centering
\repofig{fig_p3_chain_alignment.png}{0.85\linewidth}
\caption{\textbf{Contagion on the chain.} Signed alignment $A(f)$ at several
contrarian densities $\rho$, rising from the exact $f=0$ baseline
\eqref{eq:baseline} (deviation $0.0$). The axis of this run is labelled in
the per-bond convention $f_b = N_s/(N_p+2N_s)$; the per-ring values of
Eq.~\eqref{eq:f} are $f = 2 f_b$, placing the conversion onset at
$f \approx 0.67$--$0.80$.}
\label{fig:chain}
\end{figure}

\begin{figure}[t]
\centering
\repofig{fig_p3_domains_chain.png}{0.85\linewidth}
\caption{\textbf{Crossover, not a transition.} Mean $\tau$-domain length on
the chain: growth with $f$, but the $L = 8/16/32$ curves overlap with no
finite-size divergence, as expected in 1D.}
\label{fig:domains}
\end{figure}

\subsection{The conversion curve and its $\rho$-universality}

The quantity the chain already isolates is the converted fraction of
contrarians, and it collapses onto one curve in $f$, independent of $\rho$.
On the square lattice the same collapse holds and sharpens: measured at six
densities $\rho \in [0.35, 0.72]$, the conversion probability $\conv(f)$ has
maximal spread $\mathrm{std}_\rho = 0.007$, the sampling floor. The measured
curve is
%
\begin{center}
\begin{tabular}{lcccccccc}
\toprule
$f$ & $0$ & $0.5$ & $0.6$ & $0.67$ & $0.71$ & $0.75$ & $0.78$ & $0.80$ \\
$\conv(f)$ & $0$ & $0$ & $0.004$ & $0.024$ & $0.111$ & $0.278$ & $0.465$ & $0.606$ \\
\bottomrule
\end{tabular}
\end{center}
%
This is the one measured input of the theory; everything downstream is
textbook. Its $\rho$-universality is not an accident but a structural fact,
explained mechanically in the next section: a contrarian's conversion is
decided by its own witness portfolio alone.

\section{The factorization law}
\label{sec:factorization}

\subsection{Statement}

\begin{theorem}[Factorization law, prescription F]
\label{thm:factorization}
Under fresh shared witnesses, the lattice arrow field $\{\tau_i\}$ is an
independent site-percolation field: ring $i$ points forward with probability
%
\begin{equation}
\peff(\rho, f) \;=\; (1-\rho)\cdot 1 \;+\; \rho\cdot \conv(f)
\;=\; 1 - \rho\,\big(1 - \conv(f)\big),
\label{eq:peff}
\end{equation}
%
independently of every other ring. Every collective threshold follows in
closed form, with no free parameter:
%
\begin{align}
\rhoalign(f) &= \frac{1}{2\,(1-\conv(f))}
&&\text{(global reversal, } \peff = 1/2\text{)},
\label{eq:rhoalign}\\[2pt]
\rhoperc(f) &= \frac{1-p_c}{1-\conv(f)}, \quad p_c = 0.5927
&&\text{(percolation, } \peff = p_c\text{)}.
\label{eq:rhoperc}
\end{align}
\end{theorem}

The law rests on three structural facts, and its numerical proof is a
three-step falsification programme (E1--E3) that it survives.

\subsection{The three structural facts}

\emph{(1) Fresh rings never reverse (verified identity).} Under
prescription F every witness accessible to a fresh ring is fresh, and the
expected record-by-record increment of
$\langle\sigma_i^{\mathrm{loc}}\rangle$ is nonnegative, so $\tau_i = +1$
always. This rests on an exact lemma: the neighbour chiralities are
marginalized independently and the witnesses partition by bond, so the
local marginal factorizes,
$\sigma_i^{\mathrm{loc}} = \sum_{\mathrm{priv}} \sigma_p +
\sum_{\mathrm{bonds}} \sigma_b$ (checked against the engine at
$7\times 10^{-16}$), and positivity for arbitrary coordination reduces to a
single bond. The private increments are the positive closed form of
\cite{p1} (Theorem 2); the first shared increment is
$b\,\ln[(2+b)/(2-b)] \ge 0$ in closed form, with $b$ the shared record
bias; and the higher increments are positive on a $400^2$ symbolic grid
(minimum $4\times 10^{-9}$, attained only as the new record's bias
vanishes) and over $2\times 10^4$ randomized configurations of up to seven
records per bond (minimum $5\times 10^{-7}$). Exhaustively in production:
$35{,}600$ of $35{,}600$ fresh rings forward, zero backflips. The
$(1-\rho)\cdot 1$ term of Eq.~\eqref{eq:peff} is exact, and this
monotonicity is the foundation of the law: it makes every ring's fate a
local balance.

\emph{(2) Conversion depends on $f$ alone (proven).} In the true branch all
chiralities are $+1$ and contrarian status is a bath preparation local to
each witness, so the record law of ring $i$'s accessible witnesses depends
only on its own portfolio ($N_p$ charged private, $\deg \cdot N_s$ fresh
shared), never on the placement of other contrarians: the conversion
probability is a function of $f$ alone by construction. The measured
$\mathrm{std}_\rho = 0.007$ of Sec.~\ref{sec:conversion} is a check of the
engine, not evidence for the law.

\emph{(3) Conversions are independent (proven, with its only correction
localized).} $\tau_i$ is a functional of ring $i$'s accessible records
alone, records at distinct witnesses are independent given the branch, and
two rings at lattice distance $d \ge 2$ share no witness: their conversion
events are therefore exactly independent, and the arrow field is
$1$-dependent. The theory thus predicts $r(d) = 0$ identically for
$d \ge 2$, not merely small; the measured bulk correlation
$r(4 \le d \le 9) = -0.025$ is compatible (Fig.~\ref{fig:independence},
left), and the only permitted term is the nearest-neighbour correlation,
measured at $r(1) = +0.112$ ($5539$ pairs, $\sim 8\sigma$), with its
mechanism explicit: two adjacent contrarians read the same $N_s$ fresh
witnesses on their shared bond, so they convert slightly together.
% TODO(P3-F, F2): add error bars on r(2), r(3) to test the sharpened
% prediction r(d>=2) = 0 directly.

\subsection{Closed forms with no free parameter (E2)}

With $\conv(f)$ interpolated from the measured curve and no fitted parameter
anywhere, Eq.~\eqref{eq:rhoalign} reproduces the measured global-reversal
boundary with maximal residual $0.003$ (mean $0.001$), exact within
sampling, including the divergence: $\rhoalign \to \infty$ exactly when
$\conv(f) > 1/2$, i.e.\ at $f \approx 0.78$, as observed
(Sec.~\ref{sec:rhoc}). Eq.~\eqref{eq:rhoperc} reproduces the percolation
boundary with maximal residual $0.072$ (mean $0.022$); the maximum is a
single high-$f$ point, a finite-size grid crossing at $L=8$. Both overlays
are shown in Fig.~\ref{fig:factorization}. At $f = 0$,
Eq.~\eqref{eq:rhoperc} gives $\rhoperc = 1 - p_c = 0.407$, the flat
intercept of the phase map of Sec.~\ref{sec:global}, now a one-line
consequence.

\begin{figure}[t]
\centering
\repofig{fig_p3_factorization.png}{\linewidth}
\caption{\textbf{Closed-form thresholds from one measured curve (E2), no
free fit.} Left: the global-reversal boundary $\rhoalign(f)$, measured
(points) against Eq.~\eqref{eq:rhoalign} (line), maximal residual $0.003$.
Right: the percolation boundary $\rhoperc(f)$ against Eq.~\eqref{eq:rhoperc},
maximal residual $0.072$, concentrated at the highest-$f$ point (finite-size
crossing at $L=8$). $\conv(f)$ is interpolated from the measured conversion
curve; no parameter is fitted.}
\label{fig:factorization}
\end{figure}

\subsection{The ultimate test (E3)}

If Theorem~\ref{thm:factorization} is true, a purely classical, synthetic
field of independent sites at density $\peff(\rho, f)$, on the same lattices
with the same contrarian-geometry seeds, must be statistically
indistinguishable from the quantum lattice. It is: the synthetic
$\Pperc(\rho)$ curves reproduce the measured quantum ones with mean
$|\Delta| = 0.040$, below the binomial sampling floor $\approx 0.056$
(Fig.~\ref{fig:independence}, right). Within our statistics, no observable
of the arrow field distinguishes the interacting quantum lattice from
independent coin flips at density $\peff$. This is the precise, falsifiable
sense of the word \emph{reduction}.

\begin{figure}[t]
\centering
\repofig{fig_p3_independence.png}{\linewidth}
\caption{\textbf{Independence (E1) and the ultimate test (E3).} Left:
two-point correlation of contrarian conversion events versus lattice
separation $d$, bulk-compatible with zero ($r(4\le d\le 9) = -0.025$), with
a single weak nearest-neighbour correction $r(1) = +0.112$ (shared-bond
witnesses). Separations are reported for $d \le L/2$; beyond the
half-lattice the periodic images and the collapsing pair count dominate the
estimator. Right: spanning probability of the quantum arrow field (solid)
against a synthetic independent-site field at density $\peff$ on the same
geometries (dashed): mean $|\Delta| = 0.040$, below the sampling floor
$0.056$, statistically indistinguishable.}
\label{fig:independence}
\end{figure}

\subsection{The accuracy budget, and the class as a corollary}
\label{sec:budget}

The two residuals of E2 differ by a factor of twenty, and the difference is
itself a prediction of the structure. The alignment threshold
\eqref{eq:rhoalign} depends only on the mean occupation
$\langle p \rangle = \peff$: correlations between sites cancel identically
in a global average, so the $r(1)$ correction cannot touch it, and its
residual ($0.003$) sits at the sampling floor. Percolation observables, by
contrast, are sensitive to local clustering: the weak nearest-neighbour
correlation makes converted contrarians cluster slightly, shifting spanning
probabilities at the percent level, the $0.022$ mean residual of
\eqref{eq:rhoperc}. The same correlation is strictly finite-range (a single
bond), and short-range correlations are irrelevant for the critical
behaviour of percolation \cite{stauffer,weinrib}: they shift non-universal
quantities and leave the class untouched.

\begin{corollary}
The aligned-domain transition of the lattice arrow field is in the standard
two-dimensional site-percolation universality class, with $\nu = 4/3$,
$\beta/\nu = 5/48 \approx 0.104$ and $\gamma/\nu = 43/24 \approx 1.79$
\cite{stauffer,dennijs}, as a consequence of
Theorem~\ref{thm:factorization} and E3, with the $r(1)$ correction
irrelevant by \cite{weinrib}.
\end{corollary}

One microscopic mechanism, the shared bond, thus accounts for the exactness
of one threshold, the small residual of the other, and the robustness of
the class.

\section{Percolation of the global arrow}
\label{sec:global}

\subsection{A genuine transition in two dimensions}

On the square lattice the forward arrows occupy a fraction $p = (1+A)/2$ of
sites, and they can now do what the chain forbids: percolate. At
$\rho = 1/2$, the spanning probability $\Pperc(f)$ for $L \in \{4,5,6,8\}$
crosses at $f_c = 0.721$ (numerical), a finite-size percolation crossing
(Fig.~\ref{fig:percolation}); domain snapshots across the transition are
shown in Fig.~\ref{fig:snapshots}. $f_c$ sits well above the pair threshold
$f^* \approx 0.40$ for a transparent reason: in 2D the forward arrow must
reach the site-percolation threshold $p_c = 0.5927$ \cite{newmanziff}, not
merely a majority. The phase map in the $(f,\rho)$ plane
(Fig.~\ref{fig:phase}) starts at $\rho \approx 0.41$ at $f = 0$, where
$1-\rho = p_c$, and rises with sharing; its boundary is the closed form
\eqref{eq:rhoperc}, overlaid with no free parameter.

\begin{figure}[t]
\centering
\repofig{fig_p3_percolation.png}{0.85\linewidth}
\caption{\textbf{The percolation crossing.} Spanning probability
$\Pperc(f)$ at $\rho = 1/2$ for $L \in \{4,5,6,8\}$: the curves cross at
$f_c = 0.721$.}
\label{fig:percolation}
\end{figure}

\begin{figure}[t]
\centering
\repofig{fig_p3_snapshots.png}{\linewidth}
\caption{\textbf{Aligned domains across the transition.} Arrow-field
snapshots below, near, and above $f_c$: the forward domain grows from
disconnected islands to a spanning cluster.}
\label{fig:snapshots}
\end{figure}

\begin{figure}[t]
\centering
\repofig{fig_p3_2d_phase.png}{0.85\linewidth}
\caption{\textbf{Phase map of the global arrow.} Alignment $A(f,\rho)$ on
the square lattice; the measured percolation boundary (black) rises from
$\rho = 1-p_c = 0.407$ at $f=0$. Its closed form is
Eq.~\eqref{eq:rhoperc}, overlaid against the measurements in
Fig.~\ref{fig:factorization} (right).}
% TODO(P3-F, F3): _v2 of this figure adds the Eq. (7) line on the map
% itself (no free parameter); update this caption then.
\label{fig:phase}
\end{figure}

\subsection{The global arrow: shared witnesses vote double}
\label{sec:rhoc}

The critical contrarian density $\rho_c(f)$ at which the global arrow
reverses ($A = 0$) is flat at $\rho_c \approx 0.50$ up to $f \approx 0.6$,
then climbs steeply:
%
\begin{center}
\begin{tabular}{lcccccc}
\toprule
$f$ & $0.67$ & $0.71$ & $0.75$ & $0.78$ & $0.80$ & $0.83$ \\
$\rho_c$ & $0.51$ & $0.58$ & $0.70$ & $0.89$ & $>0.93$ & $>0.93$ \\
\bottomrule
\end{tabular}
\end{center}
%
Beyond $f \approx 0.78$ the forward arrow wins at every contrarian density
tested up to $0.93$: $\rho_c$ diverges. No minority, indeed no large
majority, of contrarians can reverse the global arrow. The mechanism is the
one \cite{p2} isolated in the two-ring competition: a shared witness writes
into both its rings' marginal channels at once, so its record counts
double; here the shared witnesses are fresh, so they weight the forward
arrow, and the naive majority rule $\rho_c = 1/2$ is broken upward.
Quantitatively, the entire boundary, divergence included, is the closed
form \eqref{eq:rhoalign}: the divergence point is
$\conv(f) = 1/2$, nothing more. The direction of the breaking is set by the
preparation of the shared witnesses; charged shared witnesses would break
it downward, the F/M duality again, and Sec.~\ref{sec:M} makes that
boundary sharp.

\subsection{Universality by equivalence, exponents as consistency checks}
\label{sec:universality}

The rigorous statement of universality in this paper is the Corollary of
Sec.~\ref{sec:budget}, not an exponent fit. Direct exponent estimates at
our sizes ($L \le 24$, $R = 50$ realizations; each lattice site is itself a
quantum sampling problem) then have a definite epistemic role: they are
consistency checks of the equivalence, capable of falsifying it, incapable
of establishing it, and we report them in that spirit
(Fig.~\ref{fig:convergence}; robustness panels in Appendix~\ref{app:robust}).

\emph{$\gamma/\nu$ (robust, quasi exact).} The susceptibility scaling at
the crossing gives $\gamma/\nu = 1.78$ against the percolation $1.79$,
agreement at the percent level, stable across estimators: the strongest
direct check.

\emph{$\rho_c$ (stable).} The crossing density $\rho_c^{\mathrm{perc}} =
0.407$ is stable across $L$ and matches the closed form \eqref{eq:rhoperc}.

\emph{$\nu$ (consistent, not pinned).} The slope estimator bends toward the
percolation value with increasing size (the largest pair, $L=16 \to 24$,
gives $\nu_{\mathrm{eff}} = 1.43$, drifting toward $4/3$), but $R = 50$
realizations do not pin it; collapse-metric estimates scatter in
$1.4$--$2.0$ and are $f$-dependent, a known fragility of collapse metrics
on coarse grids (the metric rewards large $\nu$, which trivializes the
rescaling). This also explains why the $f$-driven finite-size scaling of
the early runs was inconclusive: $f$ moves on the coarse integer-$N_s$
grid, and no collapse can be pinned on it; $\rho$, moving in steps of
$1/L^2$, is the native continuous control parameter. We report the drift
direction and refrain from quoting a measured $\nu$: with two usable
pairs, panel (c) of Fig.~\ref{fig:convergence} illustrates a trend, not an
extrapolation. An unequivocal direct $\nu$ requires $L \sim 64$--$128$,
i.e.\ HPC resources, and is left as the one open numerical item.

\emph{$\beta/\nu$ (weakest, as expected).} Order-parameter estimates give
$0.13$--$0.27$ against the percolation $0.104$. This is the exponent most
exposed to corrections to scaling at small $L$, and its slow convergence is
the expected finite-size pattern, not a tension with the equivalence: the
same drift toward the percolation slope is visible between the small-$L$
and large-$L$ estimates.

The consolidated claim is deliberately asymmetric: the class is established
by equivalence (E3 and the Corollary); the direct exponents are consistent
with it where they are reliable ($\gamma/\nu$), drifting toward it where
they are not yet ($\nu$, $\beta/\nu$), and nowhere against it.

\begin{figure}[t]
\centering
\repofig{fig_p3_convergence.png}{\linewidth}
\caption{\textbf{Exponent consistency checks (slope method, exact density,
$L \in \{12,16,24\}$, $R=50$).} (a) Spanning probability against exact
occupied density; the crossing gives $\rho_c = 0.407$, stable. (b) Slope
scaling at the crossing: the effective slope bends toward the percolation
value (dotted) with increasing $L$. (c) Pairwise $\nu_{\mathrm{eff}}$
against $1/L$: the drift is toward $4/3$ (dotted), shown as a trend, not a
measurement. $\gamma/\nu = 1.78$ (against percolation $1.79$) is the robust
check; $\beta/\nu$ converges slowest, as expected at these sizes.}
% TODO(P3-F, F3): switch to fig_p3_convergence_v2.png (neutral title, no
% headline exponent in title or legend).
\label{fig:convergence}
\end{figure}

\section{The boundary condition decides global time}
\label{sec:M}

Charge the shared-link baths instead of the private ones (prescription M of
\cite{p2}, $T_c = T$ in Eq.~\eqref{eq:beta}) and the factorization does not
bend; it breaks. The shared witnesses now consume records, voting backward,
and still count double. Structural fact (1) fails first and
catastrophically: $5600$ of $6400$ fresh rings reverse. The global
alignment at $\rho = 1/2$ plunges from $0$ to $-1$ as $f$ grows, monotone
down where prescription F is monotone up (Fig.~\ref{fig:prescriptionM}),
and no independent-site description survives: there is no analogue of
Eq.~\eqref{eq:peff} in which one species of ring holds a deterministic
arrow.

This is the boundary of Theorem~\ref{thm:factorization}, not a failure of
the model, and it closes the arc of the three papers. In \cite{p2} the F/M
duality decided which of two competing arrows flips; on the lattice it
decides, in addition, whether the many-arrow problem factorizes at all.
Under F the monotone structure (fresh implies forward) makes every ring's
fate a local, independent balance; under M that monotonicity is destroyed
and the collective problem is genuinely collective. As in the companion
papers, we claim the structure, not the numbers: the thresholds under M
depend on the preparation details, and what is prescription-independent is
the duality itself. The preparation of the shared records, the
Past-Hypothesis-like boundary condition that both companions deferred to
cosmology, is here promoted from deciding the sign of the arrow to deciding
the very tractability class of its emergence.

\begin{figure}[t]
\centering
\repofig{fig_p3_prescriptionM.png}{0.72\linewidth}
\caption{\textbf{The F/M duality sets the sign, and the validity of the
reduction (E4).} Global alignment $A(f)$ at $\rho = 1/2$ under the two
shared-witness preparations. F (fresh): forward wins, monotone up,
factorization holds. M (charged): backward wins, $A \to -1$, and the
factorization breaks ($5600/6400$ fresh rings reverse), so no species
carries a deterministic arrow and no independent-site reduction exists.}
\label{fig:prescriptionM}
\end{figure}

\section{Discussion}
\label{sec:discussion}

\emph{Epistemic status, consolidated} (continuing the practice of
\cite{p1,p2}). Exact: the six gates (machine precision), the $f=0$ baseline
\eqref{eq:baseline} throughout, the bond factorization of the local
marginal, structural facts (2) and (3) of the law (proven), the closed
forms \eqref{eq:rhoalign}--\eqref{eq:rhoperc} given the factorization, and
the alignment-threshold residual at the sampling floor. Verified identity:
structural fact (1), the per-bond record-increment positivity (closed form
at one record, symbolic grid and randomized configurations beyond).
Measured: the conversion curve $\conv(f)$, the one empirical input. Numerical: the
percolation phenomenology ($f_c$, the phase map), E1--E3, and the exponent
checks, with $\gamma/\nu$ quasi exact and robust, and $\nu$ convergent but
not pinned directly at these sizes; the universality claim rests on the
equivalence (E3 and the Corollary), not on the exponents.
Mechanically explained: the $\rho$-universality of $\conv$, the $r(1)$
correction (shared bond) and the twenty-fold gap between the two E2
residuals. Prescription-dependent by construction, and claimed as
structure, not numbers: the sign of the $\rho_c$ breaking and the validity
of the factorization itself (F holds, M breaks), the lattice expression of
the F/M duality.

\emph{What the reduction says.} The conjecture that closed both companion
papers is confirmed, and over-fulfilled: the knitting of local arrows is
not merely percolation-like, it is percolation, with all model physics
compressed into one measured single-ring curve. The emergence of a global
arrow of time in this model is, conditionally on its boundary conditions, a
solved problem: majority rule and site percolation of an effective forward
density. We find this compression the most striking output of the
three-paper programme: a many-body quantum problem whose collective arrow
admits a complete classical effective description with a single measured
function as interface.

\emph{What it does not say.} Everything is conditional on the preparation
of the shared witnesses, and Sec.~\ref{sec:M} shows the condition has
teeth: under the charged preparation not only the outcome but the reduction
itself is lost. The model therefore does not derive the direction of the
global arrow; it derives its propagation law given record-free shared
environments, a Past Hypothesis in miniature now doing collective work, its
origin deferred to cosmology as in \cite{p1,p2,halliwell}. If one wants a
physical image of a universally shared recording environment, the cosmic
microwave background is the natural candidate, and we leave it at that
sentence.

\emph{Limitations and openings.} The quantum lattices are small
($L \le 16$ for thresholds, $L \le 24$ for the convergence checks); the
direct $\nu$ is not pinned at these sizes, and \texttt{run\_convergence.py}
is configured for the HPC run ($L \sim 64$--$128$) that would settle it.
The factorization's structural facts are analytic: (2) and (3) are proven
from the disjointness of accessible witness sets, and (1) reduces, through
the exact bond factorization of the local marginal, to a per-bond
positivity verified symbolically and exhaustively; a closed-form proof of
that positivity for arbitrarily many shared records per bond is the one
analytic item left open. The
$\rho_c(f)$ boundary is already analytic through $\conv$
\eqref{eq:rhoalign}; the divergence region above $f \approx 0.78$ and the
exponent map at general $f$ deserve larger lattices. Natural extensions,
all starting from a closed-form baseline: other lattices and dimensions,
random graphs, and a dynamic contagion in which converted rings recharge
their neighbours' witnesses.

\appendix

\section{Definitions and the local marginal}

The lattice engine (\texttt{src/lattice/}) constructs the joint-branch
record law of \cite{p2} on arbitrary graphs: branch variable
$a_n(S) = \sum_i s_i G_{i,n}$, collision-protocol biases \eqref{eq:beta},
exact local marginals by summation over the $2^{\deg}$ neighbour
configurations. The verification suite (\texttt{verify\_lattice.py})
implements gates G1--G6 of Sec.~\ref{sec:gates} as independent recodings
(no imports from production modules), following the practice of
\cite{p1,p2}; all six pass at the residuals quoted, and the suite runs
before every production figure. The $O(2^{\deg+1})$-per-site local marginal
against the $2^{N_r}$ global object (G4, sign agreement $100\%$,
correlation $0.97$) is what makes $L=16$ (256 rings, about $2$\,s per
realization at $n_{\mathrm{traj}} = 2500$) tractable.

\section{Numerical protocols}

2D production runs: $L \in \{8,12,16\}$ (thresholds, factorization) and
$\{12,16,24\}$ (convergence checks), periodic boundary conditions,
$n_{\mathrm{traj}} = 1500$--$2500$ trajectories per site marginal, $R = 50$
disorder realizations for the convergence study; contrarian placements
resampled per realization, with geometry seeds shared between the quantum
and synthetic runs of E3 (paired comparison). Densities are reported as
exact occupied fractions $n_c/L^2$. Correlation estimates of E1 use all
contrarian pairs binned by lattice separation, reported for $d \le L/2$
(periodic-image convention); the bulk window $4 \le d \le 9$ excludes the
shared-bond correction at $d=1$. Spanning is defined by cluster wrap in
either direction. Percolation constants: $p_c = 0.5927$ (site, square
lattice) \cite{newmanziff}; exponents $\nu = 4/3$, $\beta = 5/36$,
$\gamma = 43/18$ \cite{stauffer,dennijs}. All figures regenerate from fixed
seeds by a single command; code at
\url{https://github.com/lejrimostfa/chiral-ring}.

\section{Robustness}
\label{app:robust}

\begin{figure}[h]
\centering
\repofig{fig_p3_universality.png}{\linewidth}
\caption{\textbf{Universality consistency checks at $L \le 16$.} (a)
Collapse of $\Pperc$ at the percolation $\nu = 4/3$ for two sharing
fractions; the collapse-metric optima scatter ($1.74$ at $f=0.5$, $1.38$ at
$f=0.71$), illustrating the fragility of collapse metrics at these sizes.
(b) Order parameter at the crossing: $\beta/\nu$ estimates $0.170/0.126$
against the percolation $0.104$ (dotted), the slowest-converging check.
(c) Susceptibility: $\gamma/\nu = 1.70/2.00$ bracketing the percolation
$1.79$ (dotted). All panels are consistency checks of the equivalence
(Sec.~\ref{sec:universality}), not exponent measurements.}
\label{fig:universality_fig}
\end{figure}

\section*{Acknowledgments}

This work made extensive use of Anthropic's Claude (a large language model)
as a research assistant, for model design discussions, derivations, code,
and manuscript drafting. Every analytic claim was verified by an
independently recoded numerical suite to machine precision. This
verification loop caught and corrected one convention error during
development. The author reviewed all derivations and takes sole and full
responsibility for the content.

\begin{thebibliography}{9}

\bibitem{p1} M. Lejri, \emph{Fixation of a time-reversal-odd observable by
decoherence: record distinguishability as the local arrow of time in an
exactly solvable model}, arXiv:XXXX.XXXXX (2026).

\bibitem{p2} M. Lejri, \emph{Collective irreversibility and arrow
competition in shared quantum environments}, arXiv:XXXX.XXXXX (2026).

\bibitem{zurek2009} W. H. Zurek, \emph{Quantum Darwinism}, Nat. Phys.
\textbf{5}, 181 (2009).

\bibitem{riedel2012} C. J. Riedel, W. H. Zurek, M. Zwolak, \emph{The rise
and fall of redundancy in decoherence and quantum Darwinism}, New J. Phys.
\textbf{14}, 083010 (2012).

\bibitem{newmanziff} M. E. J. Newman, R. M. Ziff, \emph{Efficient Monte
Carlo algorithm and high-precision results for percolation}, Phys. Rev.
Lett. \textbf{85}, 4104 (2000).

\bibitem{stauffer} D. Stauffer, A. Aharony, \emph{Introduction to
Percolation Theory}, 2nd ed. (Taylor \& Francis, London, 1994).

\bibitem{dennijs} M. P. M. den Nijs, \emph{A relation between the
temperature exponents of the eight-vertex and $q$-state Potts model}, J.
Phys. A \textbf{12}, 1857 (1979).

\bibitem{weinrib} A. Weinrib, B. I. Halperin, \emph{Critical phenomena in
systems with long-range-correlated quenched disorder}, Phys. Rev. B
\textbf{27}, 413 (1983).

\bibitem{halliwell} J. J. Halliwell, J. P\'erez-Mercader, W. H. Zurek
(eds.), \emph{Physical Origins of Time Asymmetry} (Cambridge Univ. Press,
1994).

\end{thebibliography}

\end{document}
